%
% einleitung.tex -- Beispiel-File für die Einleitung
%
% (c) 2020 Prof Dr Andreas Müller, Hochschule Rapperswil
%
% !TEX root = ../../buch.tex
% !TEX encoding = UTF-8
%
\section{Einleitung\label{Einleitung}}

Die Frage nach der Summe aller natürlichen Zahlen gehört zu den faszinierendsten und zugleich missverständlichsten Themen der modernen Mathematik. 
Auf den ersten Blick scheint die Reihe 
1 + 2 + 3 + 4 + \dots
offensichtlich gegen Unendlich zu divergieren. 
Dennoch taucht in verschiedenen mathematischen und physikalischen Kontexten die überraschende Aussage auf, diese Summe sei gleich 
\[
-\frac{1}{12}.
\]
Ziel dieses Papers ist es, diesen scheinbaren Widerspruch aufzulösen, indem der Begriff der Konvergenz präzisiert, die Riemannsche Zeta-Funktion eingeführt und deren analytische Fortsetzung untersucht wird. 
Abschliessend wird bewertet, in welchem Sinne die Zuordnung des Wertes
\[
-\frac{1}{12}
\] 
sinnvoll ist. 
